\chapter{Linear Transformations and Change of Basis}

In the chapters on diagonalization, the SVD, and PCA, our central move was always the same: change coordinates until the action decoupled into independent axes. Those coordinate-change formulas were practical machinery; here is the geometric principle behind them:

\begin{quote}
\emph{A linear transformation is a physical motion of space. The matrix is only
its shadow in one particular coordinate system. Change the coordinate system and
the same motion casts a different matrix.}
\end{quote}

This is the distinction between the \textbf{operator} (the geometric act of rotating, projecting, or stretching) and its \textbf{matrix} (the table of numbers written in a chosen basis). Separating the two brings clarity. Diagonalization, Jordan form, the SVD, and PCA were never isolated tricks -- they were all the same geometric sentence: \emph{pick the right basis and the machine becomes obvious.}

This chapter builds the change-of-basis engine from first principles. We define linear maps abstractly, watch how a matrix is born when a basis is chosen, establish the translation rules between different observers, and see why similarity $\mat{P}\inv\mat{A}\mat{P}$ is simply translating between two points of view. That coordinate-free perspective will prepare us directly for the complex vector spaces and Fourier analysis waiting ahead.

\section{The Operator and Its Matrix Are Two Different Things}
\label{ch31:operator-vs-matrix}

Start with a familiar geometric motion in the plane.

Take a flat sheet of paper and rotate it a quarter turn counterclockwise around
its center. That is a transformation. Every point on the paper moves; the move is
rigid (lengths and angles are preserved) and it fixes the center. Nothing about
this act mentions numbers, axes, or a grid. The rotation is just \emph{a thing
that happens to space}. Call it $T$.

Now you want to \emph{compute} with $T$. You lay a coordinate grid on the paper,
declare two unit arrows $\vec{e}_1 = \colvec{1\\0}$ and $\vec{e}_2 = \colvec{0\\1}$,
and ask: where do these two arrows go? The quarter turn sends $\vec{e}_1$ to
$\colvec{0\\1}$ and $\vec{e}_2$ to $\colvec{-1\\0}$. Stack those answers as
columns and you have the matrix
\[
  \mat{R} = \begin{bmatrix} 0 & -1 \\ 1 & 0 \end{bmatrix}.
\]
That matrix is the rotation \emph{as seen in this grid}. View the experiment in the slanted basis $\vec{b}_1=(1,1),\,\vec{b}_2=(1,-1)$, and you will get a different table of numbers for the very same physical quarter turn (Figure~\ref{ch31:fig-operator-vs-matrix}).
The rotation did not change. Your description of it did.

\begin{intuitionbox}
A linear operator is the intrinsic geometric mapping; its matrix representation depends on the chosen coordinate frame. Change the basis, and the matrix representation changes even though the operator itself remains identical. Linear algebra often seeks a basis in which the matrix representation is as simple as possible---ideally diagonal.
\end{intuitionbox}

\begin{focusbox}
Two objects, never to be confused:
\begin{itemize}
  \item A \textbf{linear operator / transformation} $T:V\to W$ -- a geometric or
        abstract rule, basis-free, the actual motion.
  \item Its \textbf{matrix} $[T]$ -- a grid of numbers that only exists \emph{after}
        you have chosen a basis for $V$ and one for $W$.
\end{itemize}
Keep this principle in mind: the matrix of a transformation always depends on
the chosen bases.
\end{focusbox}

\begin{historybox}
For two centuries following Ren\'e Descartes and Pierre de Fermat in the 1630s, geometry was dominated by Cartesian coordinates -- assigning numbers $(x,y,z)$ to points so geometric questions could be answered with algebraic equations. But this convenience came with an annoying flaw: equations depended entirely on where the axes happened to be drawn. Rotate your axes by thirty degrees, and a clean equation for an ellipse sprouted messy cross-terms like $3xy$.

In the early nineteenth century, this triggered a passionate dispute. Mathematicians who practiced \textbf{synthetic geometry} (the study of geometric figures using direct spatial relationships like intersections and tangents, without numbers) argued that algebra was blinding mathematicians to intrinsic geometric truths. Jean-Victor Poncelet complained that Cartesian formulas turned living geometry into mechanical bookkeeping. On the other side, champions of \textbf{analytic geometry} (solving geometric problems by computing with coordinate numbers) such as Arthur Cayley replied that pure spatial intuition could never scale to higher dimensions or systematic algorithms.

The conceptual tools to resolve the conflict arrived in two stages. First, in 1844, Hermann Grassmann published his theory of extension (\emph{Ausdehnungslehre}), introducing vector addition and scaling abstractly without any coordinate axes -- a vision Giuseppe Peano finally codified into modern vector-space axioms in 1888. Second, in 1872, twenty-three-year-old Felix Klein announced his famous Erlangen Program: geometry is not the study of coordinates, but the study of \textbf{invariants} -- geometric quantities that remain unchanged under a given collection of transformations.

Under this modern vision, coordinates are just a choice of basis. The real mathematical object is the coordinate-free transformation $T: V \to W$. When an observer switches from one coordinate frame to another using an invertible matrix $\mat{P}$, the matrix of the operator updates via similarity conjugation: $\mat{B} = \mat{P}\inv \mat{A} \mat{P}$. This formula is simply an interpreter translating the same underlying motion between two languages. Today, this balance remains an active frontier: in \textbf{geometric deep learning}, researchers design \textbf{equivariant} neural networks -- networks whose internal layers respect physical rotations and symmetries directly, without being fooled by arbitrary choices of coordinate frames. Meanwhile, in infinite-dimensional operator theory, the \textbf{invariant subspace problem} -- asking whether every bounded linear operator on a Hilbert space preserves at least one non-trivial closed subspace -- remains one of the most celebrated open questions in mathematics.
\end{historybox}

\begin{figure}[H]
\centering
\begin{tikzpicture}[scale=0.95,>=Stealth]
  % left panel: physical operator motion
  \begin{scope}[shift={(0,0)}]
    \node[anchor=south,font=\footnotesize\scshape,text=spacegray!90!black] at (1.5,2.6) {Physical Motion $T$ (Quarter Turn)};
    \draw[fill=codebg,draw=spacegray!30,rounded corners=4pt] (-1.8,-1.8) rectangle (3.3,2.5);
    \draw[->,gray!50] (-1.4,0) -- (2.8,0) node[right,font=\scriptsize]{$x$};
    \draw[->,gray!50] (0,-1.4) -- (0,2.1) node[above,font=\scriptsize]{$y$};
    \draw[->,vecblue,very thick] (0,0) -- (1.5,0.3) node[right,font=\small]{$\vec{x}$};
    \draw[->,vecred,very thick] (0,0) -- (-0.3,1.5) node[above,font=\small]{$T\vec{x}$};
    \draw[->,structurepurple,thick] (1.1,0.22) arc (11:101:1.12);
    \node[structurepurple,font=\scriptsize] at (1.1,1.1) {$+90^\circ$};
    \node[font=\footnotesize,align=center] at (0.75,-1.2) {rigid rotation of space\\(no coordinates needed)};
  \end{scope}

  % right panel: two matrix descriptions
  \begin{scope}[shift={(6.2,0)}]
    \node[anchor=south,font=\footnotesize\scshape,text=spacegray!90!black] at (1.8,2.9) {Two Observers, Two Matrices};
    \draw[fill=codebg,draw=spacegray!30,rounded corners=4pt] (-0.9,-2.6) rectangle (4.5,2.8);
    
    % standard basis card
    \node[draw=vecblue!60,fill=vecblue!8,rounded corners=3pt,inner sep=4pt,align=left,font=\footnotesize] at (1.8,1.5) {
      \textbf{Standard grid $\mathcal{E}$}:\\
      $\vec{e}_1 \mapsto \colvec{0\\1},\; \vec{e}_2 \mapsto \colvec{-1\\0}$\\[2pt]
      $[\mat{R}]_{\mathcal{E}} = \begin{bmatrix} 0 & -1 \\ 1 & 0 \end{bmatrix}$
    };

    % slanted basis card
    \node[draw=quotientorange!70,fill=quotientorange!8,rounded corners=3pt,inner sep=4pt,align=left,font=\footnotesize] at (1.8,-1.15) {
      \textbf{Slanted basis $\mathcal{B}=\{\vec{b}_1,\vec{b}_2\}$}:\\
      $\vec{b}_1 \mapsto \colvec{-1\\1} = -\vec{b}_2,\quad \vec{b}_2 \mapsto \colvec{1\\1} = \vec{b}_1$\\[2pt]
      $[\mat{R}]_{\mathcal{B}} = \begin{bmatrix} 0 & 1 \\ -1 & 0 \end{bmatrix}$
    };
  \end{scope}
\end{tikzpicture}
\caption{The same physical quarter turn produces different matrices in different
coordinate systems. In the basis $\vec{b}_1=(1,1),\vec{b}_2=(1,-1)$ the very same
rotation reads $\protect\begin{bsmallmatrix}0&1\\-1&0\protect\end{bsmallmatrix}$
(verified in Exercise~\ref{ch31:ex-build}). The operator is fixed; the
numbers reflect the coordinate vantage point.}
\label{ch31:fig-operator-vs-matrix}
\end{figure}

\section{Linear Maps, Abstractly}
\label{ch31:linear-maps}

Before we can change bases we need a clear statement of what a linear map even is,
in a way that does not lean on a grid at all. A \textbf{vector space} is simply any
collection of objects -- arrows in space, sequences of measurements, or polynomials --
that you can add together and scale by numbers without ever leaving the collection.
You have seen the two core rules of linearity many times; here they are as the whole
definition.

\begin{definition}[Linear transformation]
\label{ch31:def-linear}
Let $V$ and $W$ be vector spaces over $\R$. A function $T:V\to W$ is a
\textbf{linear transformation} (or \textbf{linear map}, or \textbf{linear
operator} when $W=V$) if for all $\vec{u},\vec{v}\in V$ and all scalars
$c\in\R$,
\[
  T(\vec{u}+\vec{v}) = T(\vec{u}) + T(\vec{v}),
  \qquad
  T(c\,\vec{v}) = c\,T(\vec{v}).
\]
Equivalently, $T$ preserves every linear combination:
$T(a\vec{u}+b\vec{v}) = a\,T(\vec{u}) + b\,T(\vec{v})$.
\end{definition}

That single ``preserves combinations'' line is the entire engine of the subject.
Suppose $V=\R^n$ and you know what $T$ does to
the standard basis vectors $\vec{e}_1,\dots,\vec{e}_n$. Then you know what $T$
does to \emph{everything}, because every vector is a combination of basis
vectors:
\[
  \vec{x} = x_1\vec{e}_1 + \cdots + x_n\vec{e}_n
  \;\Longrightarrow\;
  T(\vec{x}) = x_1\,T(\vec{e}_1) + \cdots + x_n\,T(\vec{e}_n).
\]
A linear map is completely determined by where it sends a basis. This is not a
small remark -- it is the reason a finite table of numbers (the matrix) can
capture a transformation of infinitely many vectors at once.

\begin{recoverybox}
\textbf{Mental check: which of these candidate maps are linear?}
Before continuing, test the two rules on three simple maps from $\R^2$ to $\R^2$:
\begin{enumerate}[label=(\alph*)]
  \item $T(x,y) = (2x - y,\, x + 3y)$
  \item $T(x,y) = (x + 1,\, y)$
  \item $T(x,y) = (x^2,\, y)$
\end{enumerate}
\emph{Answer:} Only (a) is linear. Map (b) fails because $T(0,0) = (1,0) \neq (0,0)$; shifting the origin breaks scalar scaling when $c=0$. Map (c) fails because squaring is not additive: $(u+v)^2 \neq u^2 + v^2$. A genuine linear map can only combine its input coordinates with constant multipliers.
\end{recoverybox}

\begin{aestheticbox}
Linearity leaves $T$ no freedom beyond the basis vectors: once the images of a
basis are fixed, the image of every linear combination is determined
automatically. That is why knowing $T$ on a handful of vectors pins it down
everywhere -- and why all of linear algebra reduces to ``what happens to a
basis.''
\end{aestheticbox}

\begin{example}[Two famous linear maps that are not matrices yet]
\label{ch31:ex-abstract-maps}
Linear maps arise in settings far beyond $\R^n$, which is where the abstract
definition shows its power.

\emph{Differentiation.} Let $V$ be the space of polynomials of degree at most $3$.
The derivative $D(p) = p'$ is a linear map $V\to V$:
$(\,a p + b q\,)' = a p' + b q'$. No grid in sight -- yet, as we will see, once we
pick the basis $\{1,t,t^2,t^3\}$ it \emph{becomes} a $4\times 4$ matrix.

\emph{Rotation.} The quarter turn $T$ of Section~\ref{ch31:operator-vs-matrix} is
linear: rotating a sum of arrows is the same as rotating each and adding. The
matrix $\mat{R}$ was a consequence, not the definition.
\end{example}

\begin{trapbox}
``Linear'' does not mean ``a straight-line graph'' and it does not mean
``$y=mx+b$.'' The map $x\mapsto x+1$ is \emph{not} linear: it moves the origin
($T(0)=1\neq 0$). Every genuine linear map fixes the origin, $T(\vec{0})=\vec{0}$,
because $T(\vec{0})=T(0\cdot\vec{0})=0\cdot T(\vec{0})=\vec{0}$. Shifts are
``affine,'' a strictly larger and looser family.
\end{trapbox}

\section{Coordinates: A Vector's Address in a Basis}
\label{ch31:coordinates}

To turn an abstract map into a matrix we first need to turn an abstract vector
into a column of numbers. That is what a basis does: it gives every vector an
\textbf{address}.

\begin{definition}[Coordinate vector]
\label{ch31:def-coords}
Let $\mathcal{B} = \{\vec{b}_1,\dots,\vec{b}_n\}$ be a basis of $V$. Every
$\vec{v}\in V$ is a \emph{unique} combination
$\vec{v} = c_1\vec{b}_1 + \cdots + c_n\vec{b}_n$. The list of weights is the
\textbf{coordinate vector of $\vec{v}$ relative to $\mathcal{B}$}:
\[
  [\vec{v}]_{\mathcal{B}} = \colvec{c_1\\\vdots\\c_n}.
\]
\end{definition}

Uniqueness is what makes coordinates well defined: spanning guarantees that
weights exist, and linear independence guarantees that only one tuple of weights works.
A coordinate vector is not the geometric vector itself, but its description \emph{relative to
the chosen basis $\mathcal{B}$}.

\begin{example}[The same arrow, two addresses]
\label{ch31:ex-two-addresses}
Take the geometric arrow $\vec{v}$ pointing to the point $(4,3)$ in the
plane. In the standard basis $\mathcal{E}=\{\vec{e}_1,\vec{e}_2\}$ its address is
the boring one,
\[
  [\vec{v}]_{\mathcal{E}} = \colvec{4\\3},
\]
which is why the standard basis feels invisible: the address equals the arrow.
Now switch to the slanted basis
\[
  \vec{b}_1 = \colvec{1\\1}, \qquad \vec{b}_2 = \colvec{1\\-1}.
\]
We need $c_1,c_2$ with $c_1\vec{b}_1 + c_2\vec{b}_2 = \colvec{4\\3}$, i.e.
\[
  c_1 + c_2 = 4, \qquad c_1 - c_2 = 3.
\]
Adding gives $2c_1 = 7$, so $c_1 = \tfrac72$ and $c_2 = \tfrac12$. The same arrow
now has the address
\[
  [\vec{v}]_{\mathcal{B}} = \colvec{7/2 \\ 1/2}.
\]
One arrow, two completely different lists of numbers. The arrow did not move; we
just described it to two different people who had agreed on different axes.
\end{example}

\begin{figure}[H]
\centering
\begin{tikzpicture}[scale=0.82,>=Stealth]
  % Panel A: Standard Cartesian Grid
  \begin{scope}[shift={(0,0)}]
    \node[anchor=south,font=\footnotesize\scshape,text=spacegray!90!black] at (2.2,4.3) {Standard Frame $\mathcal{E}$};
    \draw[step=1cm,gray!20,very thin] (-0.5,-1.5) grid (4.8,4.2);
    \draw[->,gray!60] (-0.5,0) -- (4.8,0) node[right,font=\scriptsize]{$x$};
    \draw[->,gray!60] (0,-1.5) -- (0,4.2) node[above,font=\scriptsize]{$y$};
    
    % standard basis vectors
    \draw[->,sourceblue,line width=1.3pt] (0,0) -- (1,0) node[below,font=\small]{$\vec{e}_1$};
    \draw[->,sourceblue,line width=1.3pt] (0,0) -- (0,1) node[above left,font=\small]{$\vec{e}_2$};
    
    % path to vector
    \draw[->,sourceblue!70,dashed,thick] (0,0) -- (4,0) node[midway,below,font=\scriptsize]{$4\vec{e}_1$};
    \draw[->,sourceblue!70,dashed,thick] (4,0) -- (4,3) node[midway,right,font=\scriptsize]{$3\vec{e}_2$};
    
    % the vector v
    \draw[->,vecred,line width=1.5pt] (0,0) -- (4,3) node[above right,font=\small]{$\vec{v}$};
    \fill[vecred] (4,3) circle (2pt);
  \end{scope}

  % Panel B: Slanted Basis Grid
  \begin{scope}[shift={(7.5,0)}]
    \node[anchor=south,font=\footnotesize\scshape,text=spacegray!90!black] at (2.2,4.3) {Slanted Frame $\mathcal{B}$};
    
    % slanted coordinate lines: lines of constant c1 (parallel to b2=(1,-1)) and c2 (parallel to b1=(1,1))
    \foreach \cone in {0,1,2,3,4} {
      \draw[gray!25,very thin] ({\cone*1-1},{\cone*1+1}) -- ({\cone*1+2},{\cone*1-2});
    }
    \foreach \ctwo in {-1,0,1,2} {
      \draw[gray!25,very thin] ({\ctwo*1-0.5},{-\ctwo*1-0.5}) -- ({\ctwo*1+4.5},{-\ctwo*1+4.5});
    }
    \draw[->,gray!60] (-0.5,0) -- (4.8,0) node[right,font=\scriptsize]{$x$};
    \draw[->,gray!60] (0,-1.5) -- (0,4.2) node[above,font=\scriptsize]{$y$};
    
    % slanted basis vectors
    \draw[->,structurepurple,line width=1.3pt] (0,0) -- (1,1) node[above left,font=\small]{$\vec{b}_1$};
    \draw[->,structurepurple,line width=1.3pt] (0,0) -- (1,-1) node[below left,font=\small]{$\vec{b}_2$};
    
    % steps along slanted basis
    \draw[->,structurepurple!80!black,dashed,thick] (0,0) -- (3.5,3.5) node[midway,above left,font=\scriptsize]{$3.5\vec{b}_1$};
    \draw[->,structurepurple!80!black,dashed,thick] (3.5,3.5) -- (4,3) node[midway,above right,font=\scriptsize]{$+0.5\vec{b}_2$};
    
    % the vector v
    \draw[->,vecred,line width=1.5pt] (0,0) -- (4,3) node[above=2pt,font=\small]{$\vec{v}$};
    \fill[vecred] (4,3) circle (2pt);
  \end{scope}
\end{tikzpicture}
\caption{One arrow, two coordinate languages. The physical vector $\vec{v}$ is identical in both panels. In the standard grid $\mathcal{E}$ (left), four steps along $\vec{e}_1$ and three steps along $\vec{e}_2$ reach $(4,3)$. In the rotated and scaled grid $\mathcal{B}$ (right), $3.5$ diagonal strides along $\vec{b}_1$ plus $0.5$ strides along $\vec{b}_2$ land on the exact same point: $[\vec{v}]_{\mathcal{B}} = (7/2,\,1/2)$.}
\label{ch31:fig-two-addresses}
\end{figure}

\begin{focusbox}
Coordinates are addresses, and addresses change when the basis changes. The
moment you write $\colvec{4\\3}$ you have silently agreed on a basis. Most of the
confusion in this chapter dissolves once you ask, every single time: \emph{a
column of numbers -- in which basis?}
\end{focusbox}

\section{The Matrix of a Transformation in Given Bases}
\label{ch31:matrix-of-T}

Now we can assemble the central object. We have an abstract linear map
$T:V\to W$. We choose a basis $\mathcal{B}=\{\vec{b}_1,\dots,\vec{b}_n\}$ for the
input space $V$ and a basis $\mathcal{C}=\{\vec{c}_1,\dots,\vec{c}_m\}$ for the
output space $W$. We want a matrix that does the same job on \emph{addresses} that
$T$ does on \emph{vectors}.

The recipe is the one we glimpsed for rotations, now stated in full generality.

\begin{formalbox}
\textbf{Matrix of a linear map.} Let $T:V\to W$ be linear, with input basis
$\mathcal{B}$ and output basis $\mathcal{C}$. For each input basis vector
$\vec{b}_j$, compute $T(\vec{b}_j)$ and write its coordinates in $\mathcal{C}$.
Those coordinate columns, side by side, form the matrix
\[
  [T]_{\mathcal{C}\leftarrow\mathcal{B}}
  = \Big[\; [T\vec{b}_1]_{\mathcal{C}} \;\big|\; [T\vec{b}_2]_{\mathcal{C}}
     \;\big|\; \cdots \;\big|\; [T\vec{b}_n]_{\mathcal{C}} \;\Big].
\]
It is the unique $m\times n$ matrix satisfying, for every $\vec{v}\in V$,
\[
  [\,T\vec{v}\,]_{\mathcal{C}} \;=\; [T]_{\mathcal{C}\leftarrow\mathcal{B}}\;[\vec{v}]_{\mathcal{B}}.
\]
\end{formalbox}

Read the subscript ``$\mathcal{C}\leftarrow\mathcal{B}$'' out loud as ``out of
$\mathcal{B}$, into $\mathcal{C}$.'' The arrow points the way the data flows: feed
in a $\mathcal{B}$-address on the right, receive a $\mathcal{C}$-address on the
left. That notation will save you from a dozen sign-and-direction errors later.

Why does the recipe work? Pure linearity. Take any $\vec{v}$ with
$\mathcal{B}$-address $\colvec{x_1\\\vdots\\x_n}$, so
$\vec{v}=\sum_j x_j\vec{b}_j$. Apply $T$ and use linearity:
\[
  T\vec{v} = \sum_{j} x_j\, T\vec{b}_j.
\]
Now take $\mathcal{C}$-coordinates of both sides. Coordinates are themselves
linear (Definition~\ref{ch31:def-coords} respects sums and scalars), so
\[
  [T\vec{v}]_{\mathcal{C}} = \sum_j x_j\,[T\vec{b}_j]_{\mathcal{C}}
  = \Big[\,[T\vec{b}_1]_{\mathcal{C}}\,\big|\cdots\big|\,[T\vec{b}_n]_{\mathcal{C}}\,\Big]\colvec{x_1\\\vdots\\x_n}.
\]
The bracketed thing is exactly the matrix above, and the column on the right is
$[\vec{v}]_{\mathcal{B}}$. Done. The columns of a transformation matrix are always
``where the input basis vectors go, written in the output basis.''

\begin{recoverybox}
\textbf{Mental check: predicting the matrix size.}
If $V$ has dimension $4$ and $W$ has dimension $3$, what are the dimensions of $[T]_{\mathcal{C}\leftarrow\mathcal{B}}$?
\emph{Answer:} $3\times 4$. The matrix has as many columns as input basis vectors ($n=4$), because we record the destination of each input basis element. It has as many rows as output basis vectors ($m=3$), because each destination is written as a list of $3$ coordinates in $\mathcal{C}$. Rows = output dimension, columns = input dimension.
\end{recoverybox}

\begin{intuitionbox}
The standard $2\times 2$ trick you already know -- ``the columns of $\mat{A}$ are
$\mat{A}\vec{e}_1$ and $\mat{A}\vec{e}_2$'' -- is just this formula with
$\mathcal{B}=\mathcal{C}=$ the standard basis. The general version says: it was
never about $\vec{e}_1,\vec{e}_2$ specifically. It is about \emph{whatever basis
you chose}. Feed the machine the basis vectors, record where they land in your
output language, stack the answers.
\end{intuitionbox}

\begin{example}[Differentiation becomes a matrix]
\label{ch31:ex-derivative-matrix}
Let $V$ be polynomials of degree $\le 3$ with basis
$\mathcal{B}=\{1,\,t,\,t^2,\,t^3\}$, and let $T=D$ be the derivative, which we
already know is linear (Example~\ref{ch31:ex-abstract-maps}). Apply $D$ to each
basis vector and express the result in the same basis $\mathcal{B}$:
\[
  D(1)=0,\quad D(t)=1,\quad D(t^2)=2t,\quad D(t^3)=3t^2.
\]
Their $\mathcal{B}$-coordinate columns are
$\colvec{0\\0\\0\\0},\colvec{1\\0\\0\\0},\colvec{0\\2\\0\\0},\colvec{0\\0\\3\\0}$,
so
\[
  [D]_{\mathcal{B}\leftarrow\mathcal{B}}
  = \begin{bmatrix} 0 & 1 & 0 & 0\\ 0 & 0 & 2 & 0\\ 0 & 0 & 0 & 3\\ 0 & 0 & 0 & 0\end{bmatrix}.
\]
Sanity check on $p(t)=5 + 2t + 7t^2 - t^3$, whose address is
$[\,p\,]_{\mathcal{B}}=\colvec{5\\2\\7\\-1}$:
\[
  [D]_{\mathcal{B}\leftarrow\mathcal{B}}\colvec{5\\2\\7\\-1}
  = \colvec{2\\14\\-3\\0},
\]
which decodes to $2 + 14t - 3t^2$. And indeed $p'(t) = 2 + 14t - 3t^2$. Calculus,
done by matrix multiplication. The abstract definition was not decoration -- it
let us treat derivatives with the same machinery as rotations.
\end{example}

\begin{connectionbox}
Treating differentiation as a
matrix (or, for infinite-dimensional function spaces, as a \emph{linear
operator}) is the founding move of \textbf{functional analysis} and of every
numerical method for differential equations. ``Discretize the derivative'' means
``replace the operator $D$ by a finite matrix in some basis,'' which is precisely
Example~\ref{ch31:ex-derivative-matrix}. Eigenvalues of $D$-like operators are
where vibration frequencies, quantum energy levels, and the heat equation's decay
rates come from.
\end{connectionbox}

\subsection{A non-square example: bases on input and output}
\label{ch31:nonsquare}

When $V$ and $W$ differ, the two bases really are independent choices and the
``out of $\mathcal{B}$, into $\mathcal{C}$'' bookkeeping becomes essential.

\begin{example}[A map from the plane to $3$-space]
\label{ch31:ex-nonsquare}
Define $T:\R^2\to\R^3$ by
$T\colvec{x\\y}=\colvec{x+y\\\,2x\,\\x-y}$. Take the standard input basis
$\mathcal{B}=\{\vec{e}_1,\vec{e}_2\}$ but a \emph{nonstandard} output basis
\[
  \mathcal{C}=\Big\{\,\vec{c}_1=\colvec{1\\0\\1},\;\vec{c}_2=\colvec{1\\0\\-1},\;\vec{c}_3=\colvec{0\\1\\0}\,\Big\}.
\]
First find where the input basis goes
(Figure~\ref{ch31:fig-nonsquare}):
$T\vec{e}_1=\colvec{1\\2\\1}$, $T\vec{e}_2=\colvec{1\\0\\-1}$. Now write each in
$\mathcal{C}$-coordinates. For $T\vec{e}_1=\colvec{1\\2\\1}$ we need
$a\vec{c}_1+b\vec{c}_2+c\vec{c}_3=\colvec{1\\2\\1}$:
\[
  a+b=1,\quad c=2,\quad a-b=1 \;\Rightarrow\; a=1,\,b=0,\,c=2,
\]
so $[T\vec{e}_1]_{\mathcal{C}}=\colvec{1\\0\\2}$. For
$T\vec{e}_2=\colvec{1\\0\\-1}=\vec{c}_2$ directly, so
$[T\vec{e}_2]_{\mathcal{C}}=\colvec{0\\1\\0}$. Hence
\[
  [T]_{\mathcal{C}\leftarrow\mathcal{B}}
  = \begin{bmatrix} 1 & 0\\ 0 & 1\\ 2 & 0 \end{bmatrix}.
\]
Same map $T$; had we used the standard output basis we would have written the
columns $\colvec{1\\2\\1},\colvec{1\\0\\-1}$ instead. The clever output basis
revealed structure the standard one hid.
\end{example}

\begin{figure}[H]
\centering
\begin{tikzpicture}[scale=1.0,>=Stealth,every node/.style={font=\small}]
  % input plane on the left
  \draw[sourceblue,thick] (-0.6,-0.8) rectangle (2.0,2.3);
  \node[sourceblue] at (0.7,2.65) {$V=\R^2$ (basis $\mathcal{B}$)};
  \draw[->,vecblue,very thick] (0,0) -- (1,0) node[below right]{$\vec{e}_1$};
  \draw[->,vecblue,very thick] (0,0) -- (0,1) node[above left]{$\vec{e}_2$};
  % output space on the right
  \node[targetgreen] at (6.4,2.65) {$W=\R^3$ (basis $\mathcal{C}$)};
  \draw[targetgreen,thick] (4.9,-0.9) rectangle (7.9,2.3);
  \node[targetgreen] at (6.4,1.4) {$[T\vec{e}_1]_{\mathcal{C}}=\colvec{1\\0\\2}$};
  \node[targetgreen] at (6.4,-0.05) {$[T\vec{e}_2]_{\mathcal{C}}=\colvec{0\\1\\0}$};
  % the arrows of the map
  \draw[->,vecred,thick] (2.1,0.8) to[bend left=12] node[midway,above]{$T$} (4.8,0.8);
\end{tikzpicture}
\caption{A non-square map carries the input-basis vectors $\vec{e}_1,\vec{e}_2$
of $V=\R^2$ into $W=\R^3$, where we record each landing spot in the
\emph{output} basis $\mathcal{C}$. Those two coordinate columns are the columns of
$[T]_{\mathcal{C}\leftarrow\mathcal{B}}$, which is therefore $3\times 2$: as many
rows as $\dimn W$, as many columns as $\dimn V$.}
\label{ch31:fig-nonsquare}
\end{figure}

\section{Change of Basis: Translating Between Two Address Systems}
\label{ch31:change-of-basis}

So far we changed the matrix of a \emph{map}. Now for the simpler-looking but
foundational task: changing the address of a fixed \emph{vector}. If two people
have agreed on two different bases, how do we translate one's address into the
other's? The answer is a single matrix, and almost everything else in the chapter
is built from it.

Suppose $\mathcal{B}=\{\vec{b}_1,\dots,\vec{b}_n\}$ and
$\mathcal{C}=\{\vec{c}_1,\dots,\vec{c}_n\}$ are two bases of the same space $V$.
We want a matrix $\mat{P}_{\mathcal{C}\leftarrow\mathcal{B}}$ that converts a
$\mathcal{B}$-address into a $\mathcal{C}$-address:
\[
  [\vec{v}]_{\mathcal{C}} = \mat{P}_{\mathcal{C}\leftarrow\mathcal{B}}\;[\vec{v}]_{\mathcal{B}}.
\]
This is just the matrix of the \emph{identity map} $\mathrm{id}:V\to V$ read out of
$\mathcal{B}$ and into $\mathcal{C}$, so the recipe of Section~\ref{ch31:matrix-of-T}
applies verbatim: \textbf{its columns are the $\mathcal{B}$ vectors written in
$\mathcal{C}$-coordinates.}
\[
  \mat{P}_{\mathcal{C}\leftarrow\mathcal{B}}
  = \Big[\,[\vec{b}_1]_{\mathcal{C}}\,\big|\,[\vec{b}_2]_{\mathcal{C}}\,\big|\cdots\big|\,[\vec{b}_n]_{\mathcal{C}}\,\Big].
\]

\begin{focusbox}
The change-of-basis matrix \emph{into} a basis is built from the \emph{other}
basis's vectors expressed in that basis. The arrow ``$\mathcal{C}\leftarrow\mathcal{B}$''
records the direction: it takes $\mathcal{B}$-coordinates as input and returns
$\mathcal{C}$-coordinates as output. Going the other way is the inverse:
$\mat{P}_{\mathcal{B}\leftarrow\mathcal{C}}=\big(\mat{P}_{\mathcal{C}\leftarrow\mathcal{B}}\big)\inv$.
\end{focusbox}

There is one case so common it deserves its own name. When $\mathcal{C}=\mathcal{E}$
is the standard basis of $\R^n$, the coordinates of $\vec{b}_j$ \emph{are just
$\vec{b}_j$ itself}. So the matrix that converts a $\mathcal{B}$-address into a
standard address is simply the basis vectors stacked as columns:
\[
  \mat{B} = \mat{P}_{\mathcal{E}\leftarrow\mathcal{B}}
  = \big[\;\vec{b}_1\;\big|\;\vec{b}_2\;\big|\cdots\big|\;\vec{b}_n\;\big].
\]
We will call this the \textbf{basis matrix} $\mat{B}$. It ``decodes'' a
$\mathcal{B}$-address back into ordinary coordinates; its inverse $\mat{B}\inv$
``encodes'' ordinary coordinates into the $\mathcal{B}$ language.

\begin{trapbox}
The basis matrix $\mat{B}$ (columns = your fancy basis vectors) goes \emph{from}
fancy coordinates \emph{to} standard ones, not the other way. People constantly
write $\mat{B}$ when they need $\mat{B}\inv$. A reliable check: $\mat{B}$ applied
to $\colvec{1\\0\\\vdots}$ must return $\vec{b}_1$ -- and it does, because that is
the first column. If your translation sends the first coordinate axis to $\vec{b}_1$,
you have the ``decode'' matrix $\mat{B}$; if it sends $\vec{b}_1$ to an axis, you
have $\mat{B}\inv$.
\end{trapbox}

\begin{example}[Translating addresses by hand]
\label{ch31:ex-translate}
Reuse the slanted basis $\vec{b}_1=\colvec{1\\1}$, $\vec{b}_2=\colvec{1\\-1}$ from
Example~\ref{ch31:ex-two-addresses}. Stack them:
\[
  \mat{B} = \begin{bmatrix} 1 & 1\\ 1 & -1\end{bmatrix},
  \qquad
  \mat{B}\inv = \frac{1}{-2}\begin{bmatrix} -1 & -1\\ -1 & 1\end{bmatrix}
  = \frac12\begin{bmatrix} 1 & 1\\ 1 & -1\end{bmatrix}.
\]
To get the $\mathcal{B}$-address of $\vec{v}=\colvec{4\\3}$ we \emph{encode} with
$\mat{B}\inv$:
\[
  [\vec{v}]_{\mathcal{B}} = \mat{B}\inv\colvec{4\\3}
  = \frac12\begin{bmatrix} 1 & 1\\ 1 & -1\end{bmatrix}\colvec{4\\3}
  = \frac12\colvec{7\\1} = \colvec{7/2\\1/2},
\]
matching the answer we found by solving equations earlier. And to go back, the
basis matrix \emph{decodes}:
\[
  \mat{B}\colvec{7/2\\1/2} = \begin{bmatrix} 1 & 1\\ 1 & -1\end{bmatrix}\colvec{7/2\\1/2}
  = \colvec{4\\3} = \vec{v}.\;\checkmark
\]
Encode with $\mat{B}\inv$, decode with $\mat{B}$. Two matrices, exact inverses,
that carry a single arrow between two languages.
\end{example}

\begin{labbox}
Do this one fully by hand on paper. Let
$\vec{b}_1=\colvec{2\\1}$, $\vec{b}_2=\colvec{1\\3}$, and $\vec{v}=\colvec{5\\5}$.
(a) Write the basis matrix $\mat{B}$. (b) Compute $\mat{B}\inv$ using the
$2\times2$ rule $\frac{1}{ad-bc}\begin{bsmallmatrix} d & -b\\ -c & a\end{bsmallmatrix}$.
(c) Find $[\vec{v}]_{\mathcal{B}}=\mat{B}\inv\vec{v}$. (d) Verify by checking that
your two weights times $\vec{b}_1,\vec{b}_2$ rebuild $\colvec{5\\5}$. You should
get weights $\colvec{2\\1}$.
\end{labbox}

\subsection{The two pictures of a basis matrix}
\label{ch31:two-pictures}

There is a subtlety here that trips up nearly everyone. The
matrix $\mat{B}$ with columns $\vec{b}_1,\vec{b}_2$ has \emph{two} jobs that look
opposite but are the same equation read two ways.

\emph{As a change of coordinates}, $\mat{B}$ takes a list of weights and decodes
it: $\mat{B}\colvec{c_1\\c_2}=c_1\vec{b}_1+c_2\vec{b}_2$. Same numbers, new
meaning.

\emph{As an active transformation}, the very same $\mat{B}$ is a machine that
\emph{moves} the standard grid to the slanted grid: it sends $\vec{e}_1\mapsto\vec{b}_1$
and $\vec{e}_2\mapsto\vec{b}_2$, dragging the unit square into the parallelogram
spanned by the basis. Whether you call $\mat{B}$ ``passive'' (relabeling the same
point) or ``active'' (moving the point) is a matter of story, not of arithmetic --
the formula $\mat{B}\vec{x}$ is identical. Hold both pictures; switch between them
deliberately.

\begin{figure}[H]
\centering
\begin{tikzpicture}[scale=0.9,>=Stealth]
  \begin{scope}
    \draw[step=1,gray!25,very thin] (0,0) grid (3,3);
    \draw[->,vecblue,very thick] (0,0) -- (1,0) node[below]{$\vec{e}_1$};
    \draw[->,vecblue,very thick] (0,0) -- (0,1) node[left]{$\vec{e}_2$};
    \fill[vecblue,opacity=0.15] (0,0)--(1,0)--(1,1)--(0,1)--cycle;
    \node at (1.5,-1) {\small standard grid};
  \end{scope}
  \draw[->,thick] (3.6,1.5) -- (5.0,1.5) node[midway,above]{$\mat{B}$};
  \begin{scope}[shift={(5.8,0)}]
    \draw[gray!25,very thin] (0,0)--(2,1); \draw[gray!25,very thin] (0,0)--(1,3);
    \draw[gray!25,very thin] (2,1)--(3,4); \draw[gray!25,very thin] (1,3)--(3,4);
    \draw[->,vecred,very thick] (0,0) -- (2,1) node[right]{$\vec{b}_1$};
    \draw[->,vecred,very thick] (0,0) -- (1,3) node[above]{$\vec{b}_2$};
    \fill[vecred,opacity=0.15] (0,0)--(2,1)--(3,4)--(1,3)--cycle;
    \node at (1.5,-1) {\small slanted grid};
  \end{scope}
\end{tikzpicture}
\caption{One matrix, two stories. $\mat{B}$ \emph{actively} carries the unit
square to the parallelogram on $\vec{b}_1,\vec{b}_2$; \emph{passively} it decodes
$\mathcal{B}$-addresses into standard ones. Same arithmetic, different narration.}
\label{ch31:fig-active-passive}
\end{figure}

\section{The Same Operator in Different Coordinates}
\label{ch31:same-operator}

Now we combine the two threads -- the matrix of a map, and translating addresses
-- into the formula this whole chapter has been walking toward. Picture an
operator $T:\R^n\to\R^n$. In the standard basis it has the familiar matrix
$\mat{A}$, so $[T\vec{v}]_{\mathcal{E}}=\mat{A}[\vec{v}]_{\mathcal{E}}$. A friend
insists on working in a basis $\mathcal{B}$ (basis matrix $\mat{B}$). What does
$T$ look like to them? We want the matrix $\mat{A}'$ with
\[
  [T\vec{v}]_{\mathcal{B}} = \mat{A}'\,[\vec{v}]_{\mathcal{B}}.
\]

Trace the action across coordinate systems to assemble the transformation matrix.
Start with a $\mathcal{B}$-address. We do not know how $T$ acts
on $\mathcal{B}$-addresses directly -- but we know how it acts on standard ones.
So:
\begin{enumerate}
  \item \textbf{Decode} the $\mathcal{B}$-address into standard coordinates with
        $\mat{B}$.
  \item \textbf{Apply $T$} in standard coordinates with $\mat{A}$.
  \item \textbf{Encode} the result back into $\mathcal{B}$-coordinates with
        $\mat{B}\inv$.
\end{enumerate}
Composing the three steps (right to left, the way matrices act):
\[
  \boxed{\;\mat{A}' = \mat{B}\inv\,\mat{A}\,\mat{B}.\;}
\]

\begin{intuitionbox}
Read $\mat{B}\inv\mat{A}\mat{B}$ as a three-step sequence:
\emph{translate into the standard language} ($\mat{B}$), \emph{apply the transformation}
($\mat{A}$), and \emph{translate back into the new language} ($\mat{B}\inv$). The two
factors of $\mat{B}$ form a round trip into and out of standard coordinates where the
action is known. The operator $T$ remains unchanged; the change-of-basis matrices
convert coordinates on the way in and out.
\end{intuitionbox}

\begin{figure}[H]
\centering
\begin{tikzpicture}[scale=0.88,>=Stealth]
  % Left panel: Commutative Diagram
  \begin{scope}[shift={(0,0)}]
    \node[anchor=south,font=\footnotesize\scshape,text=spacegray!90!black] at (2.4,3.5) {Coordinate Arithmetic};
    \draw[fill=codebg,draw=spacegray!30,rounded corners=4pt] (-1.0,-1.5) rectangle (5.8,3.4);

    % Four nodes in a square
    \node (vb) at (0,2.4) [draw=sourceblue!70,fill=sourceblue!12,rounded corners=3pt,inner sep=4pt,font=\small] {$[\vec{v}]_{\mathcal{B}}$};
    \node (wb) at (4.8,2.4) [draw=sourceblue!70,fill=sourceblue!12,rounded corners=3pt,inner sep=4pt,font=\small] {$[T\vec{v}]_{\mathcal{B}}$};
    \node (ve) at (0,0.2) [draw=targetgreen!70,fill=targetgreen!12,rounded corners=3pt,inner sep=4pt,font=\small] {$[\vec{v}]_{\mathcal{E}}$};
    \node (we) at (4.8,0.2) [draw=targetgreen!70,fill=targetgreen!12,rounded corners=3pt,inner sep=4pt,font=\small] {$[T\vec{v}]_{\mathcal{E}}$};

    % Arrows around the perimeter
    \draw[->,thick,spacegray!80!black] (vb) -- (ve) node[midway,left,font=\footnotesize,text=black]{$\mat{B}$ (decode)};
    \draw[->,thick,spacegray!80!black] (ve) -- (we) node[midway,below,font=\footnotesize,text=black]{$\mat{A}$ (transform in $\mathcal{E}$)};
    \draw[->,thick,spacegray!80!black] (we) -- (wb) node[midway,right,font=\footnotesize,text=black]{$\mat{B}\inv$ (encode)};

    % Top dashed arrow: direct shortcut
    \draw[->,vecred,line width=1.4pt,dashed] (vb) -- (wb) 
      node[midway,above=3pt,font=\footnotesize,text=vecred!90!black,draw=vecred!40,fill=white,rounded corners=2pt,inner sep=2pt] {$\mat{A}' = \mat{B}\inv\mat{A}\mat{B}$};
  \end{scope}

  % Right panel: Geometric Space (the underlying physical motion)
  \begin{scope}[shift={(8.2,0)}]
    \node[anchor=south,font=\footnotesize\scshape,text=spacegray!90!black] at (2.2,3.5) {Physical Vector Space};
    \draw[fill=codebg,draw=spacegray!30,rounded corners=4pt] (-0.8,-1.5) rectangle (5.2,3.4);

    % faint standard axes
    \draw[->,gray!40] (-0.4,0) -- (4.8,0) node[right,font=\tiny]{$x$};
    \draw[->,gray!40] (0,-0.4) -- (0,3.0) node[above,font=\tiny]{$y$};

    % physical transformation arrow from v to Tv
    \draw[->,vecblue,line width=1.5pt] (0,0) -- (2.6,0.6) node[right=2pt,font=\small]{$\vec{v}$};
    \draw[->,vecred,line width=1.5pt] (0,0) -- (1.0,2.3) node[above right=1pt,font=\small]{$T\vec{v}$};

    % physical action arc
    \draw[->,structurepurple,thick] (1.8,0.42) to[bend right=20] node[midway,above right,font=\scriptsize]{$T$} (0.8,1.8);

    % basis directions indicated
    \draw[->,structurepurple,thick,densely dotted] (0,0) -- (1.0,1.0) node[below right,font=\tiny]{$\vec{b}_1$};
    \draw[->,structurepurple,thick,densely dotted] (0,0) -- (1.2,-0.4) node[above right,font=\tiny]{$\vec{b}_2$};

    \node[font=\footnotesize,align=center] at (2.2,-1.1) {the motion $T$ is coordinate-free};
  \end{scope}
\end{tikzpicture}
\caption{The change-of-basis commutative diagram. Left: Computing the transformed address in the new basis $[T\vec{v}]_{\mathcal{B}}$ can be accomplished directly by $\mat{A}' = \mat{B}\inv\mat{A}\mat{B}$ (top dashed corridor), or equivalently by translating to standard coordinates with $\mat{B}$, executing the transformation $\mat{A}$, and encoding back with $\mat{B}\inv$. Right: The physical motion of space $T$ carrying vector $\vec{v}$ to $T\vec{v}$ is completely independent of the observer's coordinate system.}
\label{ch31:fig-commute}
\end{figure}

\begin{trapbox}
Is it $\mat{B}\inv\mat{A}\mat{B}$ or $\mat{B}\mat{A}\mat{B}\inv$? Both appear in
textbooks because they answer different questions, and getting them backwards is
the single most common change-of-basis error. The rule that never fails: the
matrix \emph{closest to the input vector on the right} must \emph{accept that
input's coordinate type}. Our input is a $\mathcal{B}$-address, and $\mat{B}$
(decode) accepts $\mathcal{B}$-addresses -- so $\mat{B}$ sits on the right. If
instead you were converting a standard-coordinate matrix \emph{from} the
$\mathcal{B}$ description, the inverse would land on the right. Always sanity-check
by asking what kind of address enters on the far right.
\end{trapbox}

\begin{example}[A shear in a different basis]
\label{ch31:ex-shear-diag}
Let $T$ have standard matrix
\[
  \mat{A} = \begin{bmatrix} 2 & 1\\ 0 & 2\end{bmatrix}.
\]
In the standard grid this is an ugly shear-plus-scale. But suppose we view it in
the basis $\vec{b}_1=\colvec{1\\1}$, $\vec{b}_2=\colvec{1\\-1}$ (basis matrix
$\mat{B}$ from Example~\ref{ch31:ex-translate}). Compute
$\mat{A}'=\mat{B}\inv\mat{A}\mat{B}$ in stages.

First $\mat{A}\mat{B}$:
\[
  \begin{bmatrix} 2 & 1\\ 0 & 2\end{bmatrix}\begin{bmatrix} 1 & 1\\ 1 & -1\end{bmatrix}
  = \begin{bmatrix} 3 & 1\\ 2 & -2\end{bmatrix}.
\]
Then $\mat{B}\inv(\mat{A}\mat{B})$ with $\mat{B}\inv=\tfrac12\begin{bsmallmatrix}1&1\\1&-1\end{bsmallmatrix}$:
\[
  \mat{A}' = \frac12\begin{bmatrix} 1 & 1\\ 1 & -1\end{bmatrix}\begin{bmatrix} 3 & 1\\ 2 & -2\end{bmatrix}
  = \frac12\begin{bmatrix} 5 & -1\\ 1 & 3\end{bmatrix}
  = \begin{bmatrix} 5/2 & -1/2\\ 1/2 & 3/2\end{bmatrix}.
\]
This particular basis did \emph{not} simplify $\mat{A}$ -- the new matrix is no
nicer. That is the real lesson: most bases make things worse. The art is finding
the special basis. For diagonalizable matrices, the natural choice is the eigenvector basis, which we turn to next.
\end{example}

\section{Similarity, Revisited}
\label{ch31:similarity}

The boxed formula $\mat{A}'=\mat{B}\inv\mat{A}\mat{B}$ is exactly the relation
called \textbf{similarity}, which we met back in Chapter~\ref{ch28:similarity} when studying the Jordan form. Here you see its geometric soul.

\begin{definition}[Similar matrices]
\label{ch31:def-similar}
Two $n\times n$ matrices $\mat{A}$ and $\mat{A}'$ are \textbf{similar} if there is
an invertible matrix $\mat{P}$ with $\mat{A}'=\mat{P}\inv\mat{A}\mat{P}$.
\end{definition}

\begin{aestheticbox}
Similarity is the statement that $\mat{A}$ and $\mat{A}'$ represent the same linear operator
expressed in different coordinate frames. The matrix $\mat{P}$ converts coordinates between the
two systems: $\mat{A}' = \mat{P}\inv\mat{A}\mat{P}$. This geometric identification explains why
algebraic invariants such as trace, determinant, and eigenvalues remain identical across all similar matrices.
\end{aestheticbox}

Because similar matrices are the same operator, they must share everything that is
truly a property of the operator and not of the chosen grid. Those shared
quantities are the \textbf{similarity invariants}, and they are precisely the
operator's coordinate-free features.

\begin{theorem}[Similar matrices share their invariants]
\label{ch31:thm-invariants}
If $\mat{A}'=\mat{P}\inv\mat{A}\mat{P}$, then $\mat{A}$ and $\mat{A}'$ have the
same determinant, the same trace, the same characteristic polynomial, and hence
the same eigenvalues (with multiplicities) and the same rank.
\end{theorem}

\begin{proof}
For the determinant, use multiplicativity:
$\det(\mat{P}\inv\mat{A}\mat{P}) = \det(\mat{P}\inv)\det(\mat{A})\det(\mat{P})
= \det(\mat{A})$, since $\det(\mat{P}\inv)\det(\mat{P})=\det(\mat{P}\inv\mat{P})=\det\mat{I}=1$.
For the characteristic polynomial, the same cancellation works with
$\mat{A}-\lambda\mat{I}$:
\[
  \mat{A}'-\lambda\mat{I} = \mat{P}\inv\mat{A}\mat{P}-\lambda\mat{P}\inv\mat{P}
  = \mat{P}\inv(\mat{A}-\lambda\mat{I})\mat{P},
\]
so $\det(\mat{A}'-\lambda\mat{I})=\det(\mat{A}-\lambda\mat{I})$ for all $\lambda$.
Equal characteristic polynomials force equal eigenvalues with multiplicities. The
trace is (up to sign) a coefficient of that polynomial -- it equals the sum of
eigenvalues -- so it is shared too; equivalently use the cyclic property
$\tr(\mat{P}\inv\mat{A}\mat{P})=\tr(\mat{A}\mat{P}\mat{P}\inv)=\tr(\mat{A})$.
Rank is preserved because multiplying by the invertible $\mat{P}$ and $\mat{P}\inv$
cannot change the dimension of the image.
\end{proof}

\begin{recoverybox}
\textbf{Mental check: can similarity change the determinant?}
Suppose a classmate claims that by choosing a clever basis $\mathcal{B}$, they turned a matrix $\mat{A}$ with $\det\mat{A} = 5$ into a similar matrix $\mat{A}'$ with $\det\mat{A}' = 1$. Could this claim possibly be true?
\emph{Answer:} Never. Similarity preserves the determinant: $\det(\mat{P}\inv\mat{A}\mat{P}) = \det\mat{A}$. Geometrically, the determinant measures the volume-scaling factor of the physical motion $T$. No change of vantage point can turn an operator that expands volume by a factor of $5$ into one that preserves volume.
\end{recoverybox}

\begin{strangbox}
Strang introduces similarity right alongside diagonalization, often through the
formula $\mat{A}=\mat{S}\Lambda\mat{S}\inv$, and emphasizes that ``$\mat{A}$ and
$\Lambda$ have the same eigenvalues.'' Reading this chapter \emph{before} Strang,
hold onto the why: eigenvalues are shared because $\mat{A}$ and $\Lambda$ are the
literally-same operator in two coordinate systems, and eigenvalues belong to the
operator, not the grid. Reading it \emph{after} Strang, notice that his $\mat{S}$
is our basis matrix $\mat{B}$ -- columns are eigenvectors -- and his
$\mat{S}\inv\mat{A}\mat{S}=\Lambda$ is our $\mat{B}\inv\mat{A}\mat{B}=\mat{A}'$
with the luckiest possible basis.
\end{strangbox}

\begin{trapbox}
Equal eigenvalues do \emph{not} imply similar. The two matrices
$\begin{bsmallmatrix}1&1\\0&1\end{bsmallmatrix}$ and
$\begin{bsmallmatrix}1&0\\0&1\end{bsmallmatrix}$ both have eigenvalue $1$ twice,
yet they are not similar: the second is the identity, and $\mat{P}\inv\mat{I}\mat{P}=\mat{I}$
for every $\mat{P}$, so the identity is similar only to itself. Sharing invariants
is necessary, not sufficient. (The full bookkeeping of what \emph{else} is needed
is the Jordan-form story.)
\end{trapbox}

\section{Choosing the Right Basis: The Natural Language of an Operator}
\label{ch31:right-basis}

We have built the machine: any basis $\mathcal{B}$ turns operator $T$ into the
matrix $\mat{B}\inv\mat{A}\mat{B}$. Most choices, as Example~\ref{ch31:ex-shear-diag}
showed, make the matrix uglier. The entire payoff of the theory is that for
special operators there is a basis that makes the matrix as simple as a matrix can
be -- and the columns of that basis matrix $\mat{B}$ are the eigenvectors or singular vectors studied across this text.

\subsection{Eigenvectors: the basis that diagonalizes}
\label{ch31:eigenbasis}

Suppose $\mat{A}$ has $n$ independent eigenvectors $\vec{v}_1,\dots,\vec{v}_n$ with
eigenvalues $\lambda_1,\dots,\lambda_n$. \emph{Use them as your basis.} Then in the
eigenvector basis the operator does the simplest thing imaginable: along each axis
$\vec{v}_j$ it just scales by $\lambda_j$ and does nothing else. The matrix
becomes diagonal. To see it from our recipe, build $\mat{B}=[\vec{v}_1|\cdots|\vec{v}_n]$
and recall $\mat{A}\vec{v}_j=\lambda_j\vec{v}_j$, so
\[
  \mat{A}\mat{B} = \big[\lambda_1\vec{v}_1\,\big|\cdots\big|\,\lambda_n\vec{v}_n\big]
  = \mat{B}\diag(\lambda_1,\dots,\lambda_n),
\]
and multiplying on the left by $\mat{B}\inv$ gives the headline result
\[
  \mat{B}\inv\mat{A}\mat{B} = \diag(\lambda_1,\dots,\lambda_n) = \Lambda.
\]
This is diagonalization, and now you see it for what it is: \emph{not} a trick, but
the statement that in the eigenvector coordinate system the operator is a pure,
axis-aligned stretch (Figure~\ref{ch31:fig-eigenstretch}).

\begin{figure}[H]
\centering
\begin{tikzpicture}[scale=0.9,>=Stealth]
  % left: standard grid, A skews the square
  \begin{scope}
    \draw[gray!30,very thin] (-0.2,-0.2) grid (2.2,2.2);
    \draw[->,vecblue,very thick] (0,0) -- (1,0);
    \draw[->,vecblue,very thick] (0,0) -- (0,1);
    % eigen-directions drawn faintly on standard grid
    \draw[->,structurepurple,thick] (0,0) -- (1.1,1.1) node[above right]{$\vec{v}_1$};
    \draw[->,quotientorange,thick] (0,0) -- (1.1,-0.55) node[below right]{$\vec{v}_2$};
    \node at (1,-1.1) {\small standard basis: $\mat{A}$ looks tangled};
  \end{scope}
  \draw[->,thick] (3.0,0.9) -- (4.4,0.9) node[midway,above]{$\mat{B}\inv\,\cdot\,\mat{B}$};
  % right: eigenbasis, pure stretch along axes
  \begin{scope}[shift={(5.4,0)}]
    \draw[gray!30,very thin] (-0.2,-0.2) grid (2.2,2.2);
    \draw[->,structurepurple,very thick] (0,0) -- (2,0) node[below right]{$\lambda_1\vec{v}_1$};
    \draw[->,quotientorange,very thick] (0,0) -- (0,1.4) node[above left]{$\lambda_2\vec{v}_2$};
    \node at (1,-1.1) {\small eigenbasis: $\Lambda$ stretches each axis};
  \end{scope}
\end{tikzpicture}
\caption{Diagonalization as a change of viewpoint. In the standard grid the
eigen-directions $\vec{v}_1,\vec{v}_2$ sit at awkward angles and $\mat{A}$ mixes
the coordinates. Rotate your axes to \emph{be} the eigenvectors and the same
operator becomes $\Lambda=\diag(\lambda_1,\lambda_2)$: each axis is simply scaled,
nothing is mixed.}
\label{ch31:fig-eigenstretch}
\end{figure}

\begin{example}[Finishing the shear example with the right basis]
\label{ch31:ex-eigenbasis}
Reconsider $\mat{A}=\begin{bsmallmatrix}2&1\\0&2\end{bsmallmatrix}$. Its only
eigenvalue is $2$ (repeated), with eigenvectors solving
$(\mat{A}-2\mat{I})\vec{v}=\vec{0}$, i.e.
$\begin{bsmallmatrix}0&1\\0&0\end{bsmallmatrix}\vec{v}=\vec0$, giving the single
direction $\vec{v}=\colvec{1\\0}$. There is only \emph{one} independent
eigenvector, not two -- so no basis can diagonalize this $\mat{A}$. It is a
\textbf{defective} matrix, and the closest we can get is the Jordan block, which is
$\mat{A}$ itself. This is the true boundary of diagonalization, and it is exactly
why the Jordan form and (for rectangular or non-diagonalizable situations) the SVD
exist.

Contrast with the symmetric matrix
$\mat{A}=\begin{bsmallmatrix}2&1\\1&2\end{bsmallmatrix}$. Its characteristic
polynomial is $(2-\lambda)^2-1=\lambda^2-4\lambda+3=(\lambda-3)(\lambda-1)$, so the
eigenvalues are $3$ and $1$, with eigenvectors $\colvec{1\\1}$ and $\colvec{1\\-1}$.
Two distinct eigenvalues force two independent eigenvectors, so the basis of those
eigenvectors makes it the clean diagonal $\diag(3,1)$. Symmetry is the friendliest
case of all -- it always yields a \emph{full, orthogonal} eigenbasis, the content
of the spectral theorem.
\end{example}

\subsection{Singular vectors: the basis that works for every matrix}
\label{ch31:svd-basis}

An eigenvector basis does not always exist: it requires a square matrix, and even
then a defective matrix comes up short (Example~\ref{ch31:ex-eigenbasis}). The singular value decomposition
fixes this by allowing \emph{two} bases, one for the input space and one for the
output, exactly the freedom we built in Section~\ref{ch31:matrix-of-T}. The SVD
$\mat{A}=\mat{U}\Sigma\mat{V}\T$ says: choose the orthonormal basis $\mathcal{V}$
(columns of $\mat{V}$) on the input side and $\mathcal{U}$ (columns of $\mat{U}$)
on the output side, and in \emph{those} coordinates $T$ is the diagonal matrix
$\Sigma$ -- it sends each input axis $\vec{v}_j$ straight to a stretched output
axis $\sigma_j\vec{u}_j$. It is the change-of-basis idea pushed to its most general
and most useful form: pick the best input basis and the best output basis
\emph{independently}, and \emph{every} linear map becomes a simple nonnegative
diagonal stretch (Figure~\ref{ch31:fig-svd}).

\begin{figure}[H]
\centering
\begin{tikzpicture}[scale=0.62,>=Stealth,every node/.style={font=\small}]
  % unit circle
  \begin{scope}
    \draw[gray!40] (0,0) circle (1);
    \draw[->,structurepurple,thick] (0,0) -- ({cos(20)},{sin(20)}) node[right]{$\vec{v}_1$};
    \draw[->,quotientorange,thick] (0,0) -- ({cos(110)},{sin(110)}) node[above]{$\vec{v}_2$};
    \node at (0,-1.7) {unit circle};
  \end{scope}
  \draw[->,thick] (1.7,0) -- (3.1,0) node[midway,above]{$\mat{V}\T$};
  % after V^T: rotate so v_i land on axes
  \begin{scope}[shift={(4.8,0)}]
    \draw[gray!40] (0,0) circle (1);
    \draw[->,structurepurple,thick] (0,0) -- (1,0);
    \draw[->,quotientorange,thick] (0,0) -- (0,1);
    \node at (0,-1.7) {axis-aligned};
  \end{scope}
  \draw[->,thick] (6.5,0) -- (7.9,0) node[midway,above]{$\Sigma$};
  % after Sigma: stretch into ellipse
  \begin{scope}[shift={(10.4,0)}]
    \draw[gray!40] (0,0) ellipse (1.6 and 0.7);
    \draw[->,structurepurple,thick] (0,0) -- (1.6,0) node[above right,inner sep=1pt]{$\sigma_1$};
    \draw[->,quotientorange,thick] (0,0) -- (0,0.7) node[above]{$\sigma_2$};
    \node at (0,-1.7) {stretched};
  \end{scope}
  \draw[->,thick] (12.4,0) -- (13.8,0) node[midway,above]{$\mat{U}$};
  % after U: rotate ellipse into output frame
  \begin{scope}[shift={(16.4,0)}]
    \draw[gray!40,rotate=25] (0,0) ellipse (1.6 and 0.7);
    \draw[->,structurepurple,thick] (0,0) -- ({1.6*cos(25)},{1.6*sin(25)}) node[above right,inner sep=1pt]{$\sigma_1\vec{u}_1$};
    \draw[->,quotientorange,thick] (0,0) -- ({-0.7*sin(25)},{0.7*cos(25)}) node[above]{$\sigma_2\vec{u}_2$};
    \node at (0,-1.7) {output};
  \end{scope}
\end{tikzpicture}
\caption{The SVD $\mat{A}=\mat{U}\Sigma\mat{V}\T$ as rotate--stretch--rotate.
$\mat{V}\T$ turns the right singular vectors onto the axes, $\Sigma$ stretches
each axis by a singular value $\sigma_j\ge 0$ (turning the circle into an
ellipse), and $\mat{U}$ rotates that ellipse into the output frame. In the
$\mathcal{V},\mathcal{U}$ bases the map is just the diagonal $\Sigma$.}
\label{ch31:fig-svd}
\end{figure}

\begin{focusbox}
The unifying slogan of the second half of the book:
\begin{center}
\emph{Diagonalization, the spectral theorem, and the SVD are all the same move --
change to the basis in which the operator is diagonal.}
\end{center}
Diagonalization uses one eigenvector basis (square, diagonalizable case). The
spectral theorem uses one \emph{orthonormal} eigenvector basis (symmetric case).
The SVD uses two orthonormal bases (every matrix, no exceptions). Each is a smarter
choice of $\mat{B}$ in $\mat{B}\inv\mat{A}\mat{B}$ -- or of $\mat{U},\mat{V}$ when
the two sides differ.
\end{focusbox}

\begin{applicationbox}
\textbf{Why data scientists obsess over the right basis.} A dataset is a cloud of
points; its raw features (``height in cm,'' ``income in dollars'') form an
arbitrary basis, usually correlated and redundant. In the previous chapter, we built
\textbf{Principal Component Analysis} to solve this exact issue. In the language of this chapter,
PCA is a change of basis: diagonalize the sample
covariance matrix (which is symmetric positive semidefinite, so the spectral
theorem guarantees an orthonormal eigenbasis) and rotate the data into that
eigenvector basis. In those new coordinates, off-diagonal covariances vanish: the
axes are completely uncoupled and ranked by variance. Keeping the top few axes
gives dimensionality reduction -- the mathematical foundation of data compression,
recommendation engines, and feature embeddings.

In deep neural networks, each layer applies a learned linear map
$\vec{x}\mapsto\mat{W}\vec{x}$ followed by a nonlinearity. Analyzing what a
layer has learned almost always comes down to inspecting its singular vectors --
the basis in which the weight matrix is a pure, uncoupled stretch. Choosing an adapted basis is central to data representation and machine learning.
\end{applicationbox}

\begin{walkbox}
Take a slow walk and turn this over: you and a friend describe the exact same
rotation of the plane, but you use the standard grid and she uses a skewed grid (say, with axes at $60^\circ$ to each other). Your matrices of numbers are different -- yet you must agree on the
angle of rotation, on the determinant ($=1$), and on the fact that there are no
real eigenvectors. \emph{Which} numbers in a matrix belong to the operator and
which belong to your grid? When you can sort every number in a matrix into one of
those two buckets, you have understood this chapter.
\end{walkbox}

\section{Putting It All Together: One Operator in Three Bases}
\label{ch31:capstone}

We end by carrying a single operator through three coordinate systems side by side.

\begin{example}[The same projection, three matrices]
\label{ch31:ex-capstone}
Let $T$ be the orthogonal projection of $\R^2$ onto the line through
$\vec{d}=\colvec{1\\1}$. Geometrically $T$ is one fixed act: flatten every arrow
onto that diagonal line. Watch its matrix representation shift as we change coordinate systems.

\emph{Standard basis.} Using the projection formula
$\proj_{\vec{d}}\vec{x}=\frac{\vec{d}\vec{d}\T}{\vec{d}\T\vec{d}}\vec{x}$ with
$\vec{d}\T\vec{d}=2$,
\[
  \mat{A} = \frac12\colvec{1\\1}\begin{bmatrix}1&1\end{bmatrix}
  = \frac12\begin{bmatrix}1&1\\1&1\end{bmatrix}
  = \begin{bmatrix}1/2&1/2\\1/2&1/2\end{bmatrix}.
\]
Off-diagonal clutter; the projection geometry is obscured by coordinate crosstalk.

\emph{Eigenvector basis.} The eigenvectors are obvious from the geometry: anything
\emph{on} the line is fixed (eigenvalue $1$) and anything \emph{perpendicular} is
killed (eigenvalue $0$). So $\vec{b}_1=\colvec{1\\1}$ (eigenvalue $1$),
$\vec{b}_2=\colvec{1\\-1}$ (eigenvalue $0$), $\mat{B}=\begin{bsmallmatrix}1&1\\1&-1\end{bsmallmatrix}$.
By the diagonalization result of Section~\ref{ch31:eigenbasis},
\[
  \mat{B}\inv\mat{A}\mat{B} = \begin{bmatrix}1&0\\0&0\end{bmatrix}.
\]
Now the matrix \emph{tells} you the geometry: keep the first axis, delete the
second. That is a projection, plainly. (Verify if you like: $\mat{A}\mat{B}
=\begin{bsmallmatrix}1&0\\1&0\end{bsmallmatrix}$, and $\mat{B}\inv$ times that is
$\begin{bsmallmatrix}1&0\\0&0\end{bsmallmatrix}$.)

\emph{Some random basis,} say $\vec{c}_1=\colvec{1\\0},\vec{c}_2=\colvec{1\\2}$,
$\mat{C}=\begin{bsmallmatrix}1&1\\0&2\end{bsmallmatrix}$,
$\mat{C}\inv=\tfrac12\begin{bsmallmatrix}2&-1\\0&1\end{bsmallmatrix}$:
\[
  \mat{C}\inv\mat{A}\mat{C}
  = \frac12\begin{bmatrix}2&-1\\0&1\end{bmatrix}\cdot\frac12\begin{bmatrix}1&1\\1&1\end{bmatrix}\cdot\begin{bmatrix}1&1\\0&2\end{bmatrix}
  = \begin{bmatrix}1/4 & 3/4\\ 1/4 & 3/4\end{bmatrix}.
\]
A third matrix for the third basis -- ugliest of all. Yet all three are similar,
so all three share $\det=0$, $\tr=1$, and eigenvalues $\{1,0\}$. Check: each has
trace $1$ ($\tfrac12+\tfrac12$, $\;1+0$, $\;\tfrac14+\tfrac34$) and determinant $0$.
The operator -- the act of flattening onto the diagonal -- never moved. We simply
described it from three vantage points, and only the eigenvector vantage made the
description transparent.
\end{example}

\begin{figure}[H]
\centering
\begin{tikzpicture}[scale=0.95,>=Stealth]
  \draw[->,gray!50] (-2,0) -- (3,0);
  \draw[->,gray!50] (0,-2) -- (0,3);
  \draw[structurepurple,thick] (-1.8,-1.8) -- (2.6,2.6) node[right]{line of $\vec{d}$};
  % a vector and its projection
  \draw[->,vecblue,very thick] (0,0) -- (2,0.4) node[right]{$\vec{x}$};
  \coordinate (p) at (1.2,1.2); % projection of (2,0.4): (2.4/2)*(1,1)=(1.2,1.2)
  \draw[->,vecred,very thick] (0,0) -- (p) node[above left]{$T\vec{x}$};
  \draw[gray!60,dashed] (2,0.4) -- (p);
  % perpendicular eigen-direction (killed)
  \draw[->,targetgreen,thick] (0,0) -- (0.9,-0.9) node[below right]{killed dir ($\lambda=0$)};
\end{tikzpicture}
\caption{The fixed operator behind Example~\ref{ch31:ex-capstone}: project onto the
line of $\vec{d}=(1,1)$. On the line, vectors are fixed ($\lambda=1$); perpendicular
to it, they are erased ($\lambda=0$). The eigenvector basis sees exactly this.}
\label{ch31:fig-capstone}
\end{figure}

\section{Summary}
\label{ch31:summary}

\begin{focusbox}
Key ideas of this chapter:
\begin{itemize}
  \item A \textbf{linear map} is basis-free; its \textbf{matrix} exists only after
        you choose bases. Columns of $[T]_{\mathcal{C}\leftarrow\mathcal{B}}$ are
        ``where the input basis vectors land, written in the output basis.''
  \item A vector's \textbf{coordinates} are its address in a basis. The
        \textbf{basis matrix} $\mat{B}$ (columns = basis vectors) decodes a
        $\mathcal{B}$-address to standard; $\mat{B}\inv$ encodes.
  \item Same operator, new basis: $\mat{A}'=\mat{B}\inv\mat{A}\mat{B}$ -- decode,
        apply operator, encode. This relation is \textbf{similarity}.
  \item Similar matrices share determinant, trace, characteristic polynomial,
        eigenvalues, and rank -- the operator's coordinate-free invariants.
  \item \textbf{The central objective} is choosing the basis that makes $\mat{A}'$
        simple: eigenvectors (diagonal), orthonormal eigenvectors (spectral
        theorem), singular vectors on both sides (SVD).
\end{itemize}
\end{focusbox}

\subsection*{The boundaries of real coordinates}

Yet as powerful as this machinery is, working strictly with real numbers and arbitrary coordinate frames runs into two fundamental boundaries.

The first boundary is geometric: \emph{many real transformations cannot be diagonalized over $\R$}. The simplest example is the quarter turn that opened this chapter. A $90^\circ$ rotation has no invariant lines in the real plane; every nonzero vector is swung into a new direction. Its characteristic polynomial is $\lambda^2 + 1 = 0$, which has no real roots. No change of real basis can ever diagonalize a rotation. Similarly, shear operators lack a full set of independent eigenvectors, leaving behind defective matrices that resist diagonal simplicity.

The second boundary is numerical: \emph{skewed basis changes can be dangerous on a computer}. If two basis vectors in $\mat{P}$ point in nearly the same direction, $\mat{P}$ becomes ill-conditioned. Computing $\mat{P}\inv\mat{A}\mat{P}$ in floating-point arithmetic then magnifies microscopic rounding errors into catastrophic noise.

Both limitations point toward the exact mathematical threshold we cross in the very next chapter.

To rescue rotations and guarantee that every operator has eigenvalues, we must enlarge our scalars from the real numbers $\R$ to the complex numbers $\C$. Over $\C$, the Fundamental Theorem of Algebra guarantees that every polynomial factors completely. Over $\C$, rotations acquire complex eigenvectors with eigenvalues $\pm i$.

To protect our computations from numerical instability, we will insist on bases whose vectors are mutually perpendicular and unit length. Over the complex numbers, these become \textbf{unitary matrices} ($\mat{U}^H\mat{U} = \mat{I}$), the complex analogues of orthogonal matrices. Unitary changes of basis never stretch lengths, never distort angles, and never amplify rounding error.

At the intersection of these two ideas stands the most celebrated unitary change of basis in modern science: the \textbf{Fourier matrix}. It maps signals from the time domain to the frequency domain, forming the basis of wave analysis, quantum mechanics, and the Fast Fourier Transform. Having mastered coordinate changes in the familiar real world, we are ready to step into the complex plane.

\section*{Exercises}
\addcontentsline{toc}{section}{Exercises}

\begin{exercise}[Quick mental check: grid versus operator]
\label{ch31:ex-quick-check}
Which of the following attributes of an $n\times n$ matrix $\mat{A}$ change when you change the basis:
(a) the individual matrix entries,
(b) the trace,
(c) the numerical coordinates of the eigenvectors,
(d) the eigenvalues,
(e) the determinant,
(f) whether the matrix is diagonalizable?
For each, state in one brief sentence whether it belongs to the grid or to the operator.
\end{exercise}

\begin{exercise}[Reading an address]
\label{ch31:ex-address}
Let $\vec{b}_1=\colvec{3\\1}$, $\vec{b}_2=\colvec{1\\2}$ and
$\vec{v}=\colvec{5\\5}$. Find $[\vec{v}]_{\mathcal{B}}$ two ways: (a) by solving
$c_1\vec{b}_1+c_2\vec{b}_2=\vec{v}$, and (b) by computing $\mat{B}\inv\vec{v}$.
Confirm the two answers agree.
\end{exercise}

\begin{exercise}[Build a transformation matrix]
\label{ch31:ex-build}
Let $T:\R^2\to\R^2$ rotate by $90^\circ$ counterclockwise. Using the basis
$\mathcal{B}=\{\colvec{1\\1},\colvec{1\\-1}\}$ for both input and output, compute
$[T]_{\mathcal{B}\leftarrow\mathcal{B}}$ directly from the recipe (find $T\vec{b}_1,
T\vec{b}_2$, then write them in $\mathcal{B}$-coordinates). Compare with
$\mat{B}\inv\mat{R}\mat{B}$ where $\mat{R}=\begin{bsmallmatrix}0&-1\\1&0\end{bsmallmatrix}$;
they must match.
\end{exercise}

\begin{exercise}[Derivative on a different basis]
\label{ch31:ex-deriv-basis}
On polynomials of degree $\le 2$, use the basis
$\mathcal{B}=\{1,\,1+t,\,1+t+t^2\}$. Find the matrix of $D=\frac{d}{dt}$ in this
basis. (Hint: compute $D$ of each basis polynomial, then express each result as a
combination of the same three basis polynomials.)
\end{exercise}

\begin{exercise}[Active versus passive]
\label{ch31:ex-active-passive}
Let $\mat{B}=\begin{bsmallmatrix}2&0\\0&1\end{bsmallmatrix}$. (a) Describe
$\mat{B}$ as an \emph{active} transformation of the unit square. (b) Describe the
\emph{same} $\mat{B}$ as a \emph{passive} decoding of $\mathcal{B}$-coordinates,
where $\mathcal{B}=\{\colvec{2\\0},\colvec{0\\1}\}$. (c) For the point with
$\mathcal{B}$-address $\colvec{1\\3}$, give its standard coordinates.
\end{exercise}

\begin{exercise}[Diagonalize by changing basis]
\label{ch31:ex-diag}
Let $\mat{A}=\begin{bsmallmatrix}4&1\\2&3\end{bsmallmatrix}$. (a) Find the
eigenvalues and eigenvectors. (b) Form the eigenvector basis matrix $\mat{B}$ and
compute $\mat{B}\inv\mat{A}\mat{B}$; verify it is diagonal with the eigenvalues on
the diagonal. (c) Use it to compute $\mat{A}^{5}$ via
$\mat{A}^5=\mat{B}\Lambda^5\mat{B}\inv$ without multiplying $\mat{A}$ by itself
five times.
\end{exercise}

\begin{exercise}[Invariants are necessary, not sufficient]
\label{ch31:ex-invariants}
(a) Show $\mat{A}=\begin{bsmallmatrix}3&0\\0&3\end{bsmallmatrix}$ and
$\mat{A}'=\begin{bsmallmatrix}3&1\\0&3\end{bsmallmatrix}$ have equal determinant,
trace, and eigenvalues. (b) Prove they are nonetheless \emph{not} similar.
(Hint: what is $\mat{P}\inv\mat{A}\mat{P}$ when $\mat{A}=3\mat{I}$, for any
invertible $\mat{P}$?)
\end{exercise}

\begin{exercise}[Trace is a similarity invariant -- prove it]
\label{ch31:ex-trace}
Prove $\tr(\mat{M}\mat{N})=\tr(\mat{N}\mat{M})$ for square matrices of compatible
size, then deduce $\tr(\mat{P}\inv\mat{A}\mat{P})=\tr(\mat{A})$. Explain in one
sentence why this means trace ``belongs to the operator,'' not the basis.
\end{exercise}

\begin{exercise}[Projection, three bases]
\label{ch31:ex-proj-three}
Let $T$ project $\R^2$ onto the line through $\colvec{2\\1}$. (a) Find the standard
matrix $\mat{A}$ using $\frac{\vec{d}\vec{d}\T}{\vec{d}\T\vec{d}}$. (b) Find the
eigenvector basis (the line direction and its perpendicular) and confirm
$\mat{B}\inv\mat{A}\mat{B}=\begin{bsmallmatrix}1&0\\0&0\end{bsmallmatrix}$.
(c) Confirm $\tr\mat{A}=1$ and $\det\mat{A}=0$ in both bases.
\end{exercise}

\begin{exercise}[Two bases, one rectangular map]
\label{ch31:ex-rect}
Let $T:\R^3\to\R^2$ be $T\colvec{x\\y\\z}=\colvec{x+z\\y-z}$. (a) Find its matrix
in the standard bases. (b) Now use the input basis
$\{\colvec{1\\0\\0},\colvec{0\\1\\0},\colvec{1\\0\\1}\}$ and the standard output
basis; recompute the matrix from the recipe and explain which column changed and
why.
\end{exercise}

\begin{exercise}[Conceptual: which numbers belong to the grid?]
\label{ch31:ex-conceptual}
A student computes the matrix of a $60^\circ$ rotation in some skewed basis and
gets entries that are not the usual $\cos,\sin$ values -- in fact the off-diagonal
entries are unequal in magnitude. Without computing, explain (i) why the entries
can look so different from the standard rotation matrix, and (ii) which two numbers
derived from the matrix she can nonetheless predict in advance, and what they are.
\end{exercise}

\begin{exercise}[Preview: the cyclic shift and roots of unity]
\label{ch31:ex-fourier-preview}
Consider the cyclic shift operator $S:\R^3\to\R^3$ defined by $S(x_1, x_2, x_3) = (x_3, x_1, x_2)$, which permutes coordinate entries in a loop.
(a) Write its matrix $\mat{P}$ in the standard basis and compute $\det\mat{P}$ and $\tr\mat{P}$.
(b) Show that the all-ones vector $\vec{v}_1 = \colvec{1\\1\\1}$ is an eigenvector with eigenvalue $\lambda_1 = 1$.
(c) Notice that repeating the shift three times returns every vector to its original state ($\mat{P}^3 = \mat{I}$). Therefore, any eigenvalue $\lambda$ must satisfy $\lambda^3 = 1$. Are there any other \emph{real} eigenvalues or eigenvectors? How does this anticipate the complex vectors and the Fourier matrix in the next chapter?
\end{exercise}
